\documentclass{article}
\usepackage[utf8]{inputenc}
\usepackage{a4wide}

\begin{document}

\section*{หาค่ากลาง}

ให้เขียนโปรแกรมเพื่อหาค่ากลาง (median) ของจำนวนเต็ม N ตัว\\

โดยที่ค่าที่รับเข้ามาจะ\textbf{เรียงลำดับมาแล้วจากน้อยไปมาก}อยู่แล้ว\\

ค่ากลางคือค่าที่อยู่ตำแหน่งตรงกลางของข้อมูล



\begin{itemize}

\item

  ถ้า N เป็นจำนวนคี่ → ค่ากลางคือค่าที่อยู่ตำแหน่งตรงกลาง

\item

  ถ้า N เป็นจำนวนคู่ → ค่ากลางคือค่าเฉลี่ยของ 2 ค่าตรงกลาง\\

  (แสดงผลเป็นทศนิยม 2 ตำแหน่งแบบปัดเศษปกติ)

\end{itemize}



\subsection*{Input}
\begin{itemize}

\item

  บรรทัดแรก: จำนวนเต็ม\texttt{N}(1 ≤ N ≤ 1,000) --- จำนวนข้อมูล

\item

  บรรทัดที่สอง: จำนวนเต็ม N ตัว\textbf{เรียงจากน้อยไปมาก}คั่นด้วยช่องว่าง

\end{itemize}



\subsection*{Output}
\begin{itemize}

\item

  แสดงค่ากลางของข้อมูล

\item

  ถ้าเป็นทศนิยม ให้แสดงทศนิยม 2 ตำแหน่ง (เช่น\texttt{25.50})

\end{itemize}



\end{document}
